\section{Cutoff fields, phase bounds, and section counts}\label{sec:C}

We continue to work with sequences of bounded dimension, and with the
subsequence convention of Section~\ref{sec:A}.  In particular, a statement
about all valuations with an order tending to zero is required on every
subsequence.  All constants below may depend on the dimension bound and
on the fixed constants explicitly present in the hypotheses.  They do
not depend on an individual member of a sequence.  Unless explicitly
stated otherwise, quantified divisorial tests are primitive prime
orders.  Homogeneous estimates extend to integral multiples; arbitrary
positive rescaling is used only for homogeneous inequalities or exact
zero conditions.

\subsection{Polarized discrepancy data and cutoff fields}
\label{subsec:C-cutoffs}

Let $V$ be a projective variety of Fano type.  Let $a$ be the discrepancy
function of an effective generalized rational lc Calabi--Yau pair over
the point, and suppose that
\[
                 a(T)\in\Z\qquad(T\text{ a primitive prime order}).
\]
Let $D$ be an effective nef and big rational $\Q$-Cartier divisor on
$V$.  Its Cartier b-pullback is denoted by $\mathbf D$.  Assume that
$b=a+\mathbf D$ is klt and has effective trace on $V$, or equivalently
$b(Q)\leq1$ for primes $Q$ on $V$.  This generalized pair is obtained
by adding $\mathbf D$ to the nef part and subtracting its trace from
the boundary; its total log canonical class is unchanged.  We call a
prime order $T$ with $a(T)=0$ a \emph{floor prime}.  Thus on $V$ every
positive coefficient of $D$ lies on the floor, and every nonfloor prime
has $a(Q)=1$ and $D(Q)=0$.

For a rational divisor $A$ on a normal projective model, put
\begin{equation}\label{eq:C-function-space}
 S(A)=\{0\}\cup\{f\in k(V)^*:\divisor(f)+A\geq0\}.
\end{equation}
Tests on this space always mean tests on its nonzero elements.  If
$A\geq0$, it contains $1$, so the field it generates is unambiguous.
We take the varying main scale $L\geq1$ to be rational.  Rational upper
approximations within a fixed factor give the same statements for real
scales, after changing constants.  A $\Q$-principal order is the order
function of a rational principal divisor.  Unless a new trace model is
specified, orders of rational divisors on $V$ mean their b-pullbacks.

We shall also use geometric generic fibres.  Divisors on such a fibre
spread over a finite extension of the generic field, and their closures
on resolutions give horizontal divisors on the total space.  Conversely,
horizontal divisors restrict divisorially.  Finite extensions of the
field of constants in characteristic zero are separable and are etale
on these generic-field models, so actual prime orders restrict to
primitive horizontal orders; integral multiples restrict to integral
multiples.  Canonical divisors from the base have horizontal order zero,
and discrepancies and Cartier b-orders therefore agree by homogeneity.
Nef b-parts restrict on a determination.  Relative numerical
proportionality for a Picard-number-one Mori contraction remains true on
a geometric generic fibre: spread any test curve over a finite extension
and specialize.  These fibres are of Fano type.  Consequently rational
numerical equivalence on them implies rational linear equivalence, by
rational connectedness and the Fano type basepoint-free consequences
recalled in Section~\ref{sec:A}.

Apply the extraction construction of Subsection~\ref{subsec:A-mori},
funded by $a+\mathbf D/2$, to obtain a projective $\Q$-factorial Fano
type model $U\to V$ containing every primitive prime with
$a=0$ and $\mathbf D\leq1$, and extracting no other prime.  The list
is finite: each selected prime has discrepancy at most $1/2$ for the
fixed generalized klt pair $a+\mathbf D/2$.  On $U$, precisely the
floor primes form $\Supp D_U$, and their coefficients lie in $(0,1]$.
For rational $0<\beta\leq1$, set
\begin{equation}\label{eq:C-cutoff}
 I_\beta=\{Q:a(Q)=0,\ D(Q)\leq\beta\},\qquad
 P_\beta=D_U-\sum_{Q\in I_\beta}D(Q)Q.
\end{equation}
All the primes in this sum occur on $U$.  Run a negative $P_\beta$-MMP
and let
\[
             p:U'\longrightarrow W_\beta
\]
be the contraction defined by the semiample output $P_\beta'$.  The
semiample divisor is the pullback of an ample rational divisor, allowing
a point as the base.  In fact there is an effective ample rational
divisor $D_W$ such that
\begin{equation}\label{eq:C-cutoff-pullback}
 P_\beta'=p^*D_W,\qquad
 D'=p^*D_W+\sum_{Q\in I_\beta}D(Q)Q'.
\end{equation}
To check the first equality as divisors, rather than just as classes,
$P_\beta'$ is effective and has no horizontal component.  A rational
principal difference between it and a pullback has divisor zero on the
projective generic fibre, hence is a function from the base.  Absorb
that function into the divisor downstairs.  Here the connected-fibre
property makes $k(W_\beta)$ relatively algebraically closed in $k(V)$.
Writing $\mathbf P$ for the nef Cartier b-divisor defined by
$P_\beta'$, MMP negativity gives
\begin{equation}\label{eq:C-cutoff-errors}
          0\leq\mathbf P\leq P_\beta^*\leq\mathbf D.
\end{equation}
Every prime in $I_\beta$ survives the MMP.  Its coefficient in
$P_\beta^*$ is zero, so the effective error
$P_\beta^*-\mathbf P$ has value zero there; a prime contracted by a
negative MMP would instead have strictly positive error.

Define
\[
                      K_\beta=k(W_\beta)\subset k(V).
\]
The field is generated by the spaces $S(mP_\beta)$, $m\in\N$, and is
relatively algebraically closed.  These fields decrease with $\beta$.
A strict inclusion of two nested relatively algebraically closed
subfields decreases transcendence degree: if their transcendence
degrees were equal, the larger would be algebraic over the smaller.
Thus there are at most $\dim V$ strict changes.  We have $K_1=k$, and
$K_{0^+}=k(V)$, since below the least coefficient in the finite support
of $D_U$ the divisor $P_\beta$ equals the big divisor $D_U$.

The following maximal-order condition, called \textup{(C-$\alpha$)}, will recur:
\begin{equation}\label{eq:C-alpha}
 T(\Gamma)\leq(1+\eta)D(T)
 \quad\text{for every effective }\Gamma\sim_\Q D,
 \qquad a(T)=0.
\end{equation}
The quantifier over $\Gamma$ is part of the condition.  When imposed on
a set of valuations in a sequence, the same $\eta=\eta_i$ works for
that entire set.

\begin{lemma}[Uniform cutoff models and order functions]
\label{lem:C-cutoff-claims}
Suppose $\dim W_\beta>0$.  After passage to a subsequence there is a
fixed $\delta>0$ and a projective $\Q$-factorial Fano type extraction
$\widehat W\to W_\beta$ such that $\widehat W$ has ordinary
$\delta$-lc singularities and every positive coefficient of the pullback
of $D_W/\beta$ is at least $\delta$.  A bounded constant $C$ can be
chosen so that $S(CD_W/\beta)$ generates $k(W_\beta)$.

There are algebraically independent functions
$\phi_1,\ldots,\phi_{\dim W_\beta}$ in this space and a b-nef
$\Q$-Cartier order function $H$ on $k(W_\beta)$, pulled back to
$k(V)$, with uniformly bounded denominator of its values on primitive
orders, such that:
\begin{enumerate}[label=\textup{(\roman*)}]
\item If $a(T)=0$, then $H(T)\leq CD(T)/\beta$.  If $T$ is nontrivial
on $K_\beta$ and all $T(\phi_j)\geq0$, then $H(T)\leq-1$.
\item If $D(Q)=0$ and $a(Q)\leq1$, then $H(Q)\leq0$.
\item If $a(T)=0$ and \eqref{eq:C-alpha} holds at $T$, then
$H(T)\geq-C\eta D(T)/\beta$.
\end{enumerate}
In particular, if \eqref{eq:C-alpha} holds on all of $I_\beta$ with
$\eta_i\to0$, those primes are eventually horizontal over $K_\beta$;
that is, their restrictions to $K_\beta$ are trivial.
\end{lemma}

\begin{proof}
On $U'\to W_\beta$ consider
\begin{equation}\label{eq:C-cutoff-direction}
 d=\frac{\mathbf D-\mathbf P}{\beta}+\mathbf P.
\end{equation}
It is nonnegative and strictly positive on $a=0$.  The discrepancies
$a+td$ have effective traces for $0\leq t\leq1$: at a surviving prime
in $I_\beta$ their trace increment is $tD(Q)/\beta\leq1$, and at an
unremoved floor prime it is $tD(Q)\leq1$.  Realize the family by a nef
increment $t\mathbf D/\beta$ and total log canonical divisor equal,
up to a rational principal divisor, to the pullback of
$t(1/\beta-1)D_W$.  Lemma~\ref{lem:A-linear-chamber} gives, on a fixed
smaller interval, discriminants
\[
                              a'+td'
\]
with $a'$ an effective generalized lc Calabi--Yau discrepancy over the
point and with uniformly bounded denominator.  These are klt for
positive $t$.

We first make explicit a lifting consequence needed here.  If
$a'(E_i)=0$, $d'(E_i)\to0$, and $k_i>0$ tends to zero, then, after a
subsequence, there are primitive lifts $T_i$ with
\begin{equation}\label{eq:C-small-lifts}
 T_i|_{k(W)}=e_i\ord_{E_i},\qquad
 a(T_i)=0,\qquad d(T_i)\to0,\qquad e_i k_i\to0.
\end{equation}
At one Mori step, choose a minimizer of $d$ on the $a$-minimum by
computing the discriminant at a sufficiently small positive parameter.
Its tailored vertical chart from Proposition~\ref{prop:A-uniform-charts}
has normalized values $0,\rho d'$ on the slice $e_*=1$.
Lemma~\ref{lem:A-lattice-approximation}, with a fixed small error
$\gamma$, gives an integral point with $e_*=lj$ bounded in terms of
$\gamma$, by rounding $ljU$ in the kernel lattice.  The error in $a$
is less than the gap of its value lattice; nonnegativity and quantization
therefore give $a=0$ exactly.  The slope bound in the same chart gives
\[
                   d\leq e_*\rho d'+C\gamma.
\]
First fix $\gamma$ and go sufficiently far in the sequence, then let
$\gamma\downarrow0$ diagonally slowly enough that also $e_*k_i\to0$.
Primitivizing the resulting order only decreases the orders and the
restriction multiplicity.  Iterate through the bounded chain of Mori
steps, at each step including the already accumulated multiplicity
times $k_i$ among the quantities required to be negligible.  Birational
steps preserve the discrepancy functions.  This proves
\eqref{eq:C-small-lifts} without asserting a uniform bound for $e_i$.

There cannot be a sequence of primes $E$ with, for one fixed sufficiently
small $t>0$,
\begin{equation}\label{eq:C-excluded-small-coefficients}
 (a'+td')(E)\longrightarrow0,
 \qquad 0<D_W(E)/\beta\longrightarrow0.
\end{equation}
The bounded denominator gives $a'(E)=0$ eventually.  Apply
\eqref{eq:C-small-lifts} with $k_i=D_W(E)/\beta$.  It gives
\[
 a(T)=0,\qquad
 \frac{D(T)}{\beta}\leq d(T)+\frac{eD_W(E)}{\beta}\longrightarrow0.
\]
Thus $T\in I_\beta$, where $\mathbf P(T)=0$, whereas
$\mathbf P(T)=eD_W(E)>0$, a contradiction.  Nor can an original prime
$E$ of $W$ have $0<D_W(E)/\beta\to0$.  A prime $Q$ above it on $U'$
satisfies
\[
 D(Q)=\mathbf P(Q)=eD_W(E)>0,\quad a(Q)=0,
 \quad a'(E)=0,\quad d'(E)\leq D_W(E).
\]
It would again give \eqref{eq:C-excluded-small-coefficients}.

Choose fixed rational $0<t\leq1/4$ in the chamber.  Extract every
exceptional prime of $A=a'+td'$ below a sufficiently small fixed
$\delta_*<1/4$.  This finite extraction is funded by mixing $a'$
slightly with an ordinary klt Calabi--Yau pair, as in
Subsection~\ref{subsec:A-mori}; the extracted primes have $a'=0$ and
positive boundary coefficient.  The preceding exclusions imply a
uniform positive lower bound on all nonzero coefficients of
\[
                 E_0=(D_W/\beta)|_{\widehat W}.
\]
Here and below the notation means the rational pullback.  The ordinary
discrepancies on the extraction dominate those of $A$, by b-nef
negativity and $\Q$-factoriality.  Every remaining exceptional primitive
prime has $A\geq\delta_*$, and every extracted prime has ordinary
discrepancy one.  This proves a uniform ordinary gap on $\widehat W$.

For completeness, we verify the additional ordinary gap needed to
apply polarized birationality to an integral Weil divisor.  Let $B_A$
be the trace boundary of $A$ on $\widehat W$.  On every component of
$E_0$ its coefficient is at least $1/2$.  For an original component the
covering-prime calculation above gives $a'=0$ and
$d'\leq D_W(\cdot)\leq1$; extracted components have $A<\delta_*$.
Choose a fixed integer $m>4/\delta$ and an integral divisor
\[
                     N=mE_0-Q_0,
\]
where the coefficients of $Q_0$ belong to $[1,2]$ on $\Supp E_0$ and
are zero elsewhere.  Such rounding is possible coefficient by
coefficient.  Since $mE_0-Q_0$ dominates a positive multiple of $E_0$,
$N$ is big.  Run an $N$-MMP to a nef and big integral Weil divisor
$N^+$ on a $\Q$-factorial model.  With stars denoting pullbacks to a
common determination, b-nef negativity and MMP negativity give
\begin{equation}\label{eq:C-rounded-errors}
 0\leq e_E=(E_0^+)^*-E_0^*,\qquad
             me_E\leq(Q_0^+)^*-Q_0^*.
\end{equation}
The total generalized log canonical divisor of $A$ is
$t(1-\beta)E_0$ up to a rational principal divisor.  Comparing its nef
part before and after the MMP, ordinary discrepancies on the output
are at least
\begin{equation}\label{eq:C-rounded-gap}
 A+(B_A^+)^*-t(1-\beta)e_E
 \ \geq\ A+\left(\frac14-\frac tm\right)(Q_0^+)^*.
\end{equation}
Indeed $B_A^+\geq Q_0^+/4$ and
$e_E\leq(Q_0^+)^*/m$ by \eqref{eq:C-rounded-errors}.
If a primitive prime with $A<\delta_*$ is in $\Supp E_0$, it has
$(Q_0^+)^*\geq Q_0^*\geq1$.  If it is not in this support, then it
is already a prime on $\widehat W$ and cannot be contracted by the
$N$-MMP: $N$ has coefficient zero there, the output pullback is
effective, and a contraction would require strictly positive MMP
error.  Its ordinary discrepancy on the output is one.  All other
primitive orders have $A\geq\delta_*$.  Formula
\eqref{eq:C-rounded-gap} therefore gives a uniform ordinary gap on the
output.  Also $N^+-K^+$ is pseudo-effective, since the output is of
Fano type.  Birkar's integral-Weil polarized birationality theorem
\cite[Theorem~1.1]{BirkarPolarised} now gives birational sections of a
bounded multiple of $N^+$.  Negativity identifies these with sections
of that multiple of $N$, hence with functions in $S(CD_W/\beta)$ for
bounded $C$.  This proves the field-generation assertion.

Choose algebraically independent $\phi_j$ in this space.  If
$\omega=\bigwedge_jd\phi_j$, define on the base
\begin{equation}\label{eq:C-logarithmic-order}
 H(E)=a'(E)-\nu_{\log,E}(\omega),\qquad
              \nu_{\log,E}(\omega)=1+\ord_E(\omega).
\end{equation}
The Calabi--Yau identity shows that $H$ is the moduli b-divisor after
a rational principal adjustment.  It is thus b-nef and $\Q$-Cartier,
and its values have bounded denominator because $a'$ does and the
log orders are integral.  If $T|_{k(W)}=e\ord_E$ and $a(T)=0$, the
discriminant minimum gives $a'(E)=0$.  Pole bounds for the $\phi_j$
on every model give
\[
                  \nu_{\log,E}(\omega)\geq-CD_W(E)/\beta.
\]
If all $E(\phi_j)\geq0$, the form is regular at that discrete
valuation ring and its logarithmic order is at least one.  These prove
(i).  For (ii), if restriction is nontrivial, then
$D_W(E)=0$ by $\mathbf P\leq\mathbf D$, so the functions are regular
and
\[
       H(Q)=e a'(E)-e\nu_{\log,E}(\omega)
                    \leq a(Q)-e\leq0.
\]
The trivial restriction gives $H(Q)=0$ directly.

On $W$ itself we also have $H_W\leq C D_W/\beta$: use $a'\leq1$,
the lower coefficient bound, and the same log-order estimate.  For any
arbitrarily small positive rational $\lambda$, choose on a determining
model an effective
\[
                  \Sigma\sim_\Q H+\lambda\mathbf D_W/\beta;
\]
this is possible because the displayed class is nef and big.  The
rational principal order
\[
           \chi=\mathbf\Sigma-H-\lambda\mathbf D_W/\beta
\]
has poles at most $C D_W/\beta$ on $W$, and therefore at most
$C\mathbf D/\beta$ upstairs.  Condition \eqref{eq:C-alpha}, applied to
$D+(\beta/C)\chi$, gives
$\chi(T)\leq C\eta D(T)/\beta$.  Since $\mathbf\Sigma\geq0$,
\[
 H(T)\geq-C\eta D(T)/\beta-\lambda\mathbf P(T)/\beta.
\]
Let $\lambda\downarrow0$ to prove (iii).  Finally, for $T\in I_\beta$,
$\mathbf P(T)=0$, so each $\phi_j$ has nonnegative order there.
If the restriction were nontrivial, (i) and (iii), together with
$D(T)\leq\beta$, would give $-1\geq H(T)\geq-C\eta_i$, which is
impossible eventually.
\end{proof}

\subsection{A phase estimate at a varying scale}

The cutoff construction identifies the part of the function field visible
at a prescribed pole cost.  We next use it to control the denominators of
rational principal orders at floor primes of very small polarization
order.  This is the phase estimate needed to regularize cone characters
in Section~\ref{sec:D}.
\label{subsec:C-scaled-phase}

\begin{lemma}[Scaled phase bound]\label{lem:C-scaled-phase}
Assume \eqref{eq:C-alpha} with $\eta_i\to0$ whenever
$a(T)=0$ and $LD(T)\to0$, along arbitrary subsequences.  Let $\ell$
be a rational principal order.  Suppose that for every nonfloor prime
$Q$ on $V$ there is an integer $p_Q$ with
\begin{equation}\label{eq:C-nonfloor-phase}
 |p_Q-\ell(Q)|\leq G(Q),\qquad
 G\geq0,\qquad G\sim_\Q c_gLD,
\end{equation}
where $c_g\geq0$ is rational and uniformly bounded.  There is a fixed
positive integer $q$ such that
\begin{equation}\label{eq:C-scaled-phase}
 \operatorname{dist}(q\ell(T),\Z)=O(LD(T))
 \quad\text{if }a(T)=0\text{ and }LD(T)\longrightarrow0.
\end{equation}
\end{lemma}

\begin{proof}
The hypothesis on nonfloor primes also holds on $U$, since every
additional extracted prime is on the floor.  If $c_g>0$, condition
\eqref{eq:C-alpha} applies to $G/(c_gL)$.  If $c_g=0$, the effective
rationally trivial divisor $G$ is zero.  Consequently, at every floor
order to which the smallness hypothesis applies,
\begin{equation}\label{eq:C-G-small}
                       G(T)\leq c_gL(1+\eta_i)D(T).
\end{equation}

Suppose the conclusion fails.  Factorial diagonalization gives a
subsequence and primitive targets $T$ with
\begin{equation}\label{eq:C-phase-failure}
 a(T)=0,\quad\beta=D(T),\quad L\beta\longrightarrow0,\qquad
 \frac{\operatorname{dist}(q\ell(T),\Z)}{L\beta}
                 \longrightarrow\infty
 \quad\text{for each fixed }q\in\N_{>0}.
\end{equation}
To justify the simultaneous last assertion, at stage $i$ choose a failure
for the multiplier $i!$ with ratio tending sufficiently quickly to
infinity.  For $q\mid i!$,
$\operatorname{dist}(i!x,\Z)\leq(i!/q)
\operatorname{dist}(qx,\Z)$; choosing the failure ratio larger than all
these factors gives the claimed diagonal sequence.
Among such sequences maximize the eventually constant
$\trdeg K_\beta$.  Lemma~\ref{lem:C-cutoff-claims} says that $T$ is
horizontal over this field.  Let $\tau$ be the first positive cutoff
for which $K_\tau=K_\beta$.  It exists because there are finitely many
coefficient cutoffs and the field $K_{0^+}=k(V)$ cannot have a nonzero
horizontal divisorial valuation.  In particular,
\begin{equation}\label{eq:C-first-change-scale}
                         0<\tau\leq\beta,
                    \qquad L\tau\longrightarrow0.
\end{equation}
After multiplying $\ell$ and $G$ by one fixed integer, maximality
implies a uniform bound
\begin{equation}\label{eq:C-before-threshold}
 \operatorname{dist}(\ell(Q),\Z)\leq CLD(Q)
                   \quad(a(Q)=0,\ D(Q)<\tau).
\end{equation}
Indeed, a failure of this assertion for every fixed multiple would,
by the same factorial diagonalization, produce a sequence at a cutoff
strictly before $\tau$, hence at strictly larger transcendence degree.

\paragraph{The first threshold.}
We first prove \eqref{eq:C-before-threshold} with $D(Q)\leq\tau$,
after another fixed multiplication.  If it fails, choose a primitive
$S$ with $D(S)=\tau$ such that every fixed multiple has phase mismatch
much larger than $L\tau$.  The prime $S$ lies on $U$.  Fix rational
$0<\gamma<1/4$ and put
\begin{equation}\label{eq:C-threshold-divisor}
 \mathcal S=S+
 \sum_{\substack{Q\text{ on }U\,,\ a(Q)=0\\Q\ne S}}
                 \gamma\min\{1,D(Q)/\tau\}\,Q.
\end{equation}
Let $t$ be the pseudo-effective threshold of $D_U-t\mathcal S$.
It is rational, since $U$ is of Fano type.  The divisor
$D_U-\tau\mathcal S$ is effective, while an effective representative
of $D_U-t\mathcal S$ would give an effective divisor rationally
equivalent to $D$ whose order at $S$ is at least $t$.
Rational pseudo-effectivity on a Fano type model is realized by such
an effective representative.  Hence \eqref{eq:C-alpha} at $S$ gives
\begin{equation}\label{eq:C-threshold-location}
                    \tau\leq t\leq(1+\eta_i)\tau,
                           \qquad\eta_i\longrightarrow0.
\end{equation}
Run the threshold MMP and take the semiample nonbig fibration.  On its
generic fibre,
\begin{equation}\label{eq:C-threshold-comparison}
 D'\sim_\Q t\mathcal S',\qquad
 \sum_{Q\ne S}D(Q)Q'\precsim2\eta_i\tau S'.
\end{equation}
Here $A\precsim B$ means that $B-A$ is rationally linearly equivalent
to an effective rational divisor.  To see the second inequality, the
coefficient of every $Q'\ne S'$ in $D'-t\mathcal S'$ is at least
$D(Q)/2$ for large $i$, whereas the coefficient of $S'$ is at least
$-\eta_i\tau$.  Restriction of the rationally trivial threshold
class proves the claim.  Enlarge $\eta_i$ rationally if necessary.
The restriction of $D'$ is big, so $S'$ is horizontal and big on this
generic fibre.  For the threshold MMP and every subsequent MMP over
its base, the pullback errors satisfy
\begin{equation}\label{eq:C-threshold-error}
 0\leq e_D=(D'')^*-\mathbf D
       \leq t\bigl((\mathcal S'')^*-\mathcal S^*\bigr).
\end{equation}
The first comparison uses that $\mathbf D$ is nef; the second is
MMP negativity for $D-t\mathcal S$, followed by its crepancy over the
threshold base.

We need a signed integral multiplier, allowed to grow slowly, for which
\begin{equation}\label{eq:C-phase-amplification}
 m=m_i\ne0,\quad\sigma=\{m\ell(S)\}\longrightarrow0^+,
                    \qquad\eta_i+|m|L\tau=o(\sigma).
\end{equation}
Pass first to a limiting phase.  If it is irrational, choose fixed
signed powers whose limiting residues approach zero from above, and
then let these powers vary sufficiently slowly.  For a rational limit,
first take a fixed multiple clearing that limit.  Its signed residual
$\zeta_i\to0$ satisfies $|\zeta_i|/(L\tau)\to\infty$ by the assumed
failure.  Choose its sign positive and multiply it by an integer so
that the resulting residual still tends to zero but dominates
$\eta_i$.  This is possible with no wrap: choose a target residual
$s_i\to0$ with $s_i\gg\eta_i+|\zeta_i|$, and a multiplier of size
$s_i/|\zeta_i|$.  Its product with $L\tau$, divided by the new
residual, is $O(L\tau/|\zeta_i|)\to0$.  This proves
\eqref{eq:C-phase-amplification}.

On the threshold model choose a rational correction $n$ so that
$N=m\ell'+n$ is an integral Weil divisor, as follows:
\begin{align}
 n(S')&=-\sigma,\label{eq:C-rounded-phase-data}\\
 |n(Q')|&\leq C|m|LD(Q)
                 &&(Q\ne S,\ a(Q)=0,\ D(Q)<\tau),\notag\\
 0\leq n(Q')&<1
                 &&(Q\ne S,\ a(Q)=0,\ D(Q)\geq\tau),\notag\\
 |n(Q')|&\leq |m|G'(Q')
                 &&(a(Q)>0).\notag
\end{align}
Use nearest integers in the second line, upward rounding in the third,
and $mp_Q$ in the last; outside the finite relevant support the
correction is zero.  Since $\ell'$ is rational principal,
\eqref{eq:C-threshold-comparison} bounds the positive errors in class
by
\[
 C|m|LD'+|m|G'+\tau^{-1}\sum_{Q\ne S}D(Q)Q'
          \precsim O(|m|L\tau+\eta_i)S'.
\]
Thus
\[
                    N\precsim-(\sigma-o(\sigma))S',
\]
and $N$ is not relatively pseudo-effective.  Run its Mori fibre MMP
over the threshold base.  On the geometric generic fibre $F$ of the
last Picard-number-one contraction,
\begin{equation}\label{eq:C-large-index-threshold}
 H_1=-N''\text{ is integral and ample},\qquad
                   H_1\precsim(\sigma+o(\sigma))S''.
\end{equation}
The upper comparison uses the lower bounds in
\eqref{eq:C-rounded-phase-data}; all the generic class comparisons
survive pushing and restricting.  The variety $F$ is of Fano type,
and $-K_F$ dominates $S''|_F$ numerically.  In fact the pushed crepant
boundary of $a$ is effective and contains $S''$ with coefficient one,
while the moduli trace on $F$ is pseudo-effective.  For this last
assertion restrict its nef determination first, add arbitrarily small
ample errors and a pullback of an ample divisor if needed, and push
effective approximations to $F$.  Letting the errors tend to zero
proves pseudo-effectivity.  Picard-number-one proportionality on the
total contraction, followed by its geometric generic restriction and
rational linear equivalence on $F$, therefore gives
\begin{equation}\label{eq:C-unbounded-index}
                         -K_F\sim_\Q I_iH_1,
                               \qquad I_i\longrightarrow+\infty.
\end{equation}

We verify the small-discrepancy hypothesis that contradicts
Lemma~\ref{lem:B-unbounded-index}.  Let $w$ be an integral divisorial
angular evaluation on $F$, allowing a nonprimitive one, with ordinary
discrepancy $A_F(w)\to0$.  Generalized-to-ordinary comparison gives
\begin{equation}\label{eq:C-small-angular}
                  A_F(w)\geq a(w)+w((\mathcal S'')^*).
\end{equation}
Indeed the pushed boundary has coefficient one at every surviving
floor prime, dominates $\mathcal S''$, and b-nef negativity controls
the moduli trace.  Hence integrality gives $a(w)=0$ eventually and,
by \eqref{eq:C-threshold-error},
$w(\mathcal S^*)\leq w((\mathcal S'')^*)\to0$.
The formula on $U$ gives $\mathbf D\leq C\mathcal S^*$, with fixed
$C$ because $D(Q)\leq1$ and $\tau\leq1$.  Thus the primitive
restriction of $w$ is eventually one of the primes extracted on $U$.
Since $w(\mathcal S^*)\to0$, that prime is neither $S$ nor a floor
prime of coefficient at least $\tau$: their $\mathcal S$-coefficients
are respectively $1$ and $\gamma$.  The remaining coefficient is
$\gamma D(Q)/\tau$, so
\begin{equation}\label{eq:C-small-D-threshold}
                   D(w)\leq C\tau A_F(w).
\end{equation}
The solved estimate \eqref{eq:C-before-threshold}, extended by integral
homogeneity, and \eqref{eq:C-phase-amplification} imply that
$m\ell(w)$ has a residual $o(A_F(w))$ modulo $\Z$.  Equations
\eqref{eq:C-G-small} and \eqref{eq:C-threshold-error} also give
\begin{equation}\label{eq:C-G-threshold-error}
 |m|w((G'')^*)
       =|m|\bigl(G(w)+c_gLe_D(w)\bigr)=o(A_F(w)).
\end{equation}
The equality holds because $G-c_gLD$ is rational principal and its
pullback is unchanged.  In the estimate, $e_D(w)\leq
 t\,w((\mathcal S'')^*)=O(\tau A_F(w))$.

The signs of the rounding corrections now matter.  Their negative
part is at most
\[
 \sigma(S'')^*+C|m|L(D'')^*+|m|(G'')^*;
\]
its value is $o(A_F(w))$ by the preceding bounds.  The positive
unit-size errors are supported on floor components, whose reduced
sum has order at most $A_F(w)$ by boundary comparison.  Thus
\begin{equation}\label{eq:C-signed-rounding}
                   -o(A_F(w))\leq w((n'')^*)\leq O(A_F(w)).
\end{equation}
All these are inequalities between effective divisors pulled back from
$\Q$-factorial models before restriction to $F$.

Consider the compact Seifert cone associated to the integral
polarization $H_1$.  For a divisorial evaluation over infinity with
positive height $u$ and ordinary discrepancy $A\to0$,
Lemma~\ref{lem:B-seifert} gives $A\geq u+A_F(w)$ and
\[
                    u\equiv-m\ell(w)-w((n'')^*)\pmod\Z.
\]
The trivial angular restriction is impossible when $A\to0$, since
then $u$ is a positive integer.  Otherwise all terms tend to zero,
so there is no wrap in this congruence.  The residual of $m\ell(w)$
is $o(A_F(w))$ and the negative part in
\eqref{eq:C-signed-rounding} is $o(A_F(w))$; positivity of $u$
therefore forces $u=o(A)$.  Orders with centre outside infinity have $v(D)=0$, so the same
small-height conclusion holds for every small-discrepancy prime.
This contradicts Lemma~\ref{lem:B-unbounded-index} and
\eqref{eq:C-unbounded-index}.  We have proved the phase estimate also
at $D(S)=\tau$, after a fixed multiplication.

\paragraph{Rational triviality on the cutoff fibre.}
Return to the cutoff model $U^\tau\to W_\tau$.  On its generic
fibre,
\[
                   D'=\sum_{Q\in I_\tau}D(Q)Q',
\]
and no other floor prime on this model is horizontal.  With the fixed
multiplications already made, choose an integral Weil divisor
$N=\ell'+n$ whose horizontal trace satisfies
\begin{equation}\label{eq:C-cutoff-rounding}
                         |n|\leq CLD'+G'.
\end{equation}
If either $N$ or $-N$ is not relatively pseudo-effective, run its Mori
fibre MMP.  On the final geometric generic Mori fibre, put
$R=\sum_{Q\in I_\tau}Q''$.  The same pullback and class comparisons
give
\begin{equation}\label{eq:C-second-large-index}
 H_1=\mp N''\precsim O(L\tau)R,\qquad
          0\leq\mathbf D\leq(D'')^*\leq\tau R^*.
\end{equation}
Again $-K$ dominates $R$ numerically, so the integral ample divisor
$H_1$ has anticanonical index tending to infinity.  If $A_F(w)\to0$,
ordinary comparison gives $a(w)=0$, $w(R^*)\leq A_F(w)$, and hence
\[
 D(w)=O(\tau A_F(w)),\qquad
                  w((G'')^*)=O(L\tau A_F(w)).
\]
For the latter estimate use \eqref{eq:C-G-small} at the primitive
restriction and the positive trace-pullback error bounded by
$(D'')^*$.  That primitive restriction is eventually in $I_\tau$,
since $D(w)\to0$, and all such primes survive on $U^\tau$.
Its solved phase gives a residual $o(A_F(w))$ for $\ell(w)$.
Equation \eqref{eq:C-cutoff-rounding}, pulled back after the MMP,
gives $w((n'')^*)=o(A_F(w))$.  The Seifert congruence consequently
forces $u=o(A)$ as before, contradicting
Lemma~\ref{lem:B-unbounded-index}.

Both $N$ and $-N$ are therefore relatively pseudo-effective.  On a
Fano type generic fibre they each have effective rationally equivalent
representatives, by semiample termination.  Their sum is an effective
rationally trivial divisor and is zero, so $N$ is rationally linearly
trivial on that fibre.

\paragraph{A bounded principalization multiple.}
Run an MMP over $W_\tau$ making $R=\sum_{I_\tau}Q''$ nef.  On the
geometric generic fibre the integral divisor $N''$ remains rationally
principal, and its pullback agrees with that of $N$, since a
rationally trivial divisor is crepant through these modifications.
Fix $0<c<1$.  The discrepancies
\begin{equation}\label{eq:C-torsion-CY}
                              a_*=a+cR^*
\end{equation}
are an effective generalized klt Calabi--Yau pair.  The trace is
effective because $R$ is reduced and lies on the floor; klt follows
from $\mathbf D\leq\tau R^*$ and positivity of $\mathbf D$ on
$a=0$.  The additional b-part is nef.  The same small-order argument
shows
\[
       \operatorname{dist}(N''(w),\Z)=o(a_*(w))
                          \quad\text{whenever }a_*(w)\to0.
\]
We record why this bounds its torsion order.  If its minimal
principalization order $l$ were unbounded, write
$lN''=\divisor(g)$ and take the normalization in a root of $g$.
The cover is connected by minimality of $l$ and is unramified in
codimension one because $N''$ is integral.  Thus ordinary Fano type
data pull back without an angular ramification boundary.  Form the
product with $\mathbf P^1$ and take the finite diagonal cyclic
quotient, as in Lemma~\ref{lem:B-principal-order}.  The product-to-quotient
map is quasi-etale; over zero the fibration has multiplicity $l$.
The log pullback of \eqref{eq:C-torsion-CY}, product with the line
factor, and descent provide effective generalized klt Calabi--Yau
data over the curve.

For a prime over zero, put
$u=\ord(\text{base coordinate})/l>0$.  Its discrepancy satisfies
$A\geq u+a_*(w)$, and the cyclic congruence is
$u\equiv-w(N'')\pmod\Z$.  A trivial angular restriction makes $u$
a positive integer and hence cannot occur when $A\to0$.  For a
nontrivial restriction the preceding phase estimate and the absence of
wrap give $u=o(A)$.  Add the rational divisor $f^*[0]/l$ with varying coefficient, extending
the height by zero away from the origin.  The discrepancies $A-su$ remain an
effective generalized Calabi--Yau family over the curve for arbitrarily
long nonnegative parameter intervals.  Lemma~\ref{lem:A-comparison}
applies because every $A\to0$ has $u=o(A)$.  But the base discrepancy
at zero is at most $1/l\to0$ and is nonpositive after the coefficient
$s=1$ is added.  This contradicts its conclusion.  Thus $l$ is bounded.

The bounded principalization descends to the original generic field.
Cartierness descends under the algebraically closed field extension for
the corresponding Weil divisor, and a geometrically trivial line
bundle and its inverse each have a nonzero section over the original
field by flat base change.  Their product is a nonzero constant, so
those sections give a trivialization.  This also explains explicitly
why no unbounded field extension is absorbed into the asserted bounded
multiple.

\paragraph{The final two-sided comparison.}
There is now a fixed positive integer $q$ and a rational function
$\chi$ on the generic fibre, hence in $k(V)$, such that the rational
principal order
\[
                         h=q\ell-\divisor(\chi)
\]
has horizontal trace on $U^\tau$ bounded by
\begin{equation}\label{eq:C-final-horizontal-bound}
                            |h'|\leq CLD'+qG'.
\end{equation}
For example, take $\divisor(\chi)=qN$ on that fibre, so
$h'=-qn$.  The original target $T$ is horizontal over $K_\tau=K_\beta$,
and subtracting an integral principal order does not remove its phase
failure:
\begin{equation}\label{eq:C-final-h-failure}
                          \frac{|h(T)|}{LD(T)}\longrightarrow\infty.
\end{equation}
On the geometric generic fibre consider
\begin{equation}\label{eq:C-final-family}
 a_0=a+c\mathbf D/\tau,\qquad
 a_s=a_0+\frac{s h-q|s|\mathbf G}{J_iL\tau}.
\end{equation}
Choose $J_i\to\infty$ sufficiently slowly compared with the ratio in
\eqref{eq:C-final-h-failure}, and then a two-sided parameter range
$|s|\leq R_i\to\infty$ with $R_i/J_i\to0$.  These are concave,
finite piecewise affine effective generalized Calabi--Yau data.
Indeed \eqref{eq:C-final-horizontal-bound} and
$D'\leq\tau\sum_{I_\tau}Q'$ give trace effectivity for this range;
off the floor the term $-q|s|\mathbf G$ absorbs the possible positive
trace of $sh$.  Up to rational principal terms the added nef part has
$\mathbf D$-coefficient
$c/\tau-q|s|c_g/(J_i\tau)>0$.  Thus b-nefness also holds.

The central member is klt.  If a primitive geometric generic order
satisfies
$a_0<\min\{1,c\}/2$, then $a=0$ and $D<\tau/2$.
Its primitive horizontal restriction is therefore a surviving prime
of $I_\tau$ on $U^\tau$.  Equations
\eqref{eq:C-G-small} and \eqref{eq:C-final-horizontal-bound} give
\[
                        |h|+q\mathbf G=O(L\mathbf D)
\]
there, uniformly for all such orders: these are primes already on the
model where the horizontal trace bound was proved.  Thus the fixed
discrepancy cutoff controls every order required by the two-sided
comparison lemma.  For every fixed signed $s$ the
relative difference $(a_s-a_0)/a_0$ is consequently $O(|s|/J_i)\to0$.
This is the strong two-sided hypothesis of
Lemma~\ref{lem:A-comparison}, with a fixed positive discrepancy
cutoff.  Its conclusion over a point says that $a_s$ is klt for fixed
$s$ eventually.  At the lift of $T$, however, choose $s=1$ or $s=-1$
so that $sh(T)=-|h(T)|$.  Equations
\eqref{eq:C-final-family} and \eqref{eq:C-final-h-failure}, with the
choice of $J_i$, give $a_s(T)<0$.  This contradiction proves the lemma.
\end{proof}

\subsection{Exact saturation of neutral fields}

For the pinning argument, a small error on a fixed list of functions is
not enough: later steps use rational expressions of unrestricted degree
in those functions.  The next lemma puts every bounded-cost function in
one relatively algebraically closed cutoff field and makes the relevant
valuations exactly trivial on that field.  Its count estimates use one
normalization for both the upper and lower bounds.
\label{subsec:C-saturation}

\begin{lemma}[Saturation and section counts]\label{lem:C-saturation-count}
Let $f_1,\ldots,f_s$ be a finite list of genuine rational functions on
$V$, whose length may vary with the sequence, and let
\begin{equation}\label{eq:C-neutral-polarization}
 J(T)=\max\{0,\max_j(-T(f_j))\},\qquad
 0\leq J\leq C L_0\mathbf D,\qquad1\leq L_0=O(L).
\end{equation}
Thus $J$ is an integral basepoint-free Cartier b-divisor on a high
model.  Suppose \eqref{eq:C-alpha} holds uniformly with $\eta_i\to0$
on
\[
                       \mathcal L=\{T:a(T)=0,\ J(T)=0\}.
\]
In addition, on every subsequence on which $L_0\to\infty$, assume
that all valuations in $\mathcal L$ eventually restrict trivially to
the \emph{field generated by} $S(L_0D)$.

After passage to a subsequence, there are cutoff fields $K_*$ of fixed
transcendence degree $p$, a fixed constant $B$, and positive real
numbers $v_i$ with the following properties.
\begin{enumerate}[label=\textup{(\roman*)}]
\item For every fixed $C_1>0$, eventually $S(C_1LD)\subset K_*$.
\item Eventually $S(BLD)$ generates $K_*$.  If $T\in\mathcal L$ is
nonnegative on this space, its restriction to $K_*$ is trivial.
\item There is $C'>0$, independent of the positive integer $n$, such
that for every fixed $n$, eventually
\begin{equation}\label{eq:C-uniform-counts}
 \dim S(nLD)\leq C'v_i n^p,\qquad
 \dim\operatorname{Span}\{g_1\cdots g_n:g_j\in S(BLD)\}
                                      \geq v_i n^p.
\end{equation}
\end{enumerate}
The eventual index in \textup{(i)} may depend on $C_1$, and that in
\textup{(iii)} may depend on $n$; $B$ and $C'$ are fixed before $n$.
All assertions concern divisorial evaluations.  Exact triviality also
holds on rational monomial subcones cut out by the indicated
conditions.
\end{lemma}

\begin{proof}[Proof of the field assertions]
Because $J$ is integral and $L_0=O(L)$, choose a fixed small $c>0$ so
that
\[
                  a(T)=0,\quad LD(T)\leq c
                            \quad\Longrightarrow\quad J(T)=0.
\]
There is, after a subsequence, a single field plateau containing two
cutoffs
\begin{equation}\label{eq:C-plateau}
 0<\beta_0\leq\beta_1\leq c/L,\qquad
 L\beta_0\longrightarrow0,\qquad L\beta_1\geq c'>0,\qquad
                         K_{\beta_0}=K_{\beta_1}=K_*.
\end{equation}
Here is a precise reason for this choice.  Pass to limits of the finitely
many change locations multiplied by $L$, adding dummy endpoints if the
number of changes varies.  The last change tending to zero lies below
a cutoff $\beta_0$ with $L\beta_0\to0$; choose it decreasing slowly
enough to lie beyond that change.  The next change either has a positive
limiting scaled location or lies beyond $c/L$.  Choose $\beta_1$
strictly before it, with $L\beta_1$ fixed positive and at most $c$.
This gives \eqref{eq:C-plateau}.  Passing again fixes $p=\trdeg K_*$.

All $I_{\beta_1}$ belong to $\mathcal L$, so
Lemma~\ref{lem:C-cutoff-claims} makes them horizontal over $K_*$.
It follows, for \emph{every} rational $m>0$, that
\begin{equation}\label{eq:C-exact-plateau-spaces}
                   S(mP_{\beta_0})=S(mP_{\beta_1}).
\end{equation}
Indeed a function in the left side belongs to $K_*$: enlarge $m$ to
an integer and use the section-field description.  Its order at every
removed prime in $I_{\beta_1}$ is zero by horizontality.  Removing
those pole allowances therefore leaves it in the right side.  The
reverse inclusion follows from $P_{\beta_1}\leq P_{\beta_0}$.
For fixed $C_1$, a function in $S(C_1LD)$ has no pole at
$I_{\beta_0}$ eventually, because its order is integral and
$C_1L\beta_0<1$.  Hence
\[
          S(C_1LD)\subset S(C_1LP_{\beta_0})
                       =S(C_1LP_{\beta_1})\subset K_*.
\]
This proves (i).  The cutoff generation statement at $\beta_1$, and
$\beta_1\geq c'/L$, place a birational field-generating space for
$K_*$ inside $S(B_0LD)$ for fixed $B_0$.  Enlarging $B$ will therefore
make $S(BLD)$ generate $K_*$; a basis contains $p$ algebraically
independent elements, by greedily choosing a transcendence basis from
any field-generating set.

We prove that some fixed $B$ forces the asserted exact triviality.
Otherwise choose $B_i\to\infty$ slowly and $T_i\in\mathcal L$ that
are nonnegative on $S(B_iLD)$ but nontrivial on $K_*$.  Since $K_1=k$,
there is a change at $\beta\geq\beta_1$ where $T$ is trivial on
$K_\beta$ but nontrivial on the preceding field $K_{\beta_-}$.
Take $\beta_-$ sufficiently close to $\beta$ from below that
$\beta_->\beta/2$.  Use $H,\phi_j$ from
Lemma~\ref{lem:C-cutoff-claims} at $\beta_-$.  Their pole bounds are
$C\mathbf D/\beta$, and $1/\beta=O(L)$; thus $T(\phi_j)\geq0$
eventually.  Its nontrivial restriction gives
\begin{equation}\label{eq:C-negative-cutoff-order}
                                  H(T)\leq-1.
\end{equation}

If $L_0\beta\to\infty$, then $L_0\to\infty$ because $\beta\leq1$.
The same bounded-cost field-generating space for $K_{\beta_-}$ is
contained in $S(L_0D)$ eventually.  The hypothesis of triviality on the
\emph{field} generated by the latter space contradicts the choice of
$T$.  It remains to treat a subsequence on which $L_0\beta$ is
bounded.

Restrict to the geometric generic fibre of the cutoff contraction
$U^\beta\to W_\beta$.  There $T$ is horizontal and
\[
                     D'\leq\beta\sum_{I_\beta}Q'.
\]
Apply the extra-data part of Lemma~\ref{lem:A-linear-chamber}, over a
point, with
\begin{equation}\label{eq:C-saturation-chamber}
 \widetilde a=a+\kappa J,\qquad
             d=c_* =\mathbf D/\beta,\qquad h=H,
\end{equation}
where $\kappa>0$ is fixed rational and sufficiently small; $c_*$
denotes the additional nonnegative direction of that lemma.
We verify all its hypotheses.  The orders $\widetilde a$ have bounded
denominator because $a$ and $J$ are integral.  The order $H$ has bounded
denominator by the cutoff lemma.  Both $J$ and $\mathbf D/\beta$ are
b-nef, and $H$ is pulled back from a b-nef divisor on $K_{\beta_-}$;
after a birational determination it restricts to a b-nef divisor on
this fibre.  Their Calabi--Yau realizations hold the total class
rationally trivial.  At floor primes on this fibre,
\[
                    H\leq C D/\beta,\qquad J\leq C D/\beta,
\]
where the second bound uses boundedness of $L_0\beta$.  At other
primes on the model, $a=1$, $D=0$, $J=0$, and $H\leq0$.
These inequalities give effective traces on uniform small positive
ranges in \eqref{eq:C-saturation-chamber}, including the family with
$\lambda H+sc_*$ for $0<s\leq\lambda\theta_0$.
The family $\widetilde a+td$ is klt for $t>0$ because $D$ is positive
on $a=0$.  Finally,
\[
 \widetilde a=0\quad\Longleftrightarrow\quad a=J=0,
 \qquad -C\eta_i c_*\leq H\leq Cc_*
                       \quad\text{on this zero set}.
\]
The lower bound is part (iii) of the cutoff lemma, using the assumed
uniform \eqref{eq:C-alpha}.  Thus the exact sign conclusion of
Lemma~\ref{lem:A-linear-chamber} over a point prohibits $H<0$ at the
lift of $T$, contradicting \eqref{eq:C-negative-cutoff-order}.
This proves (ii).

The proof of (iii), including its common normalization $v_i$, is given
below after establishing the adjoint lattice-point statement.  For the
extension to rational monomial subcones, take any finite list of tested
functions, resolve their divisors, and subdivide the cone so their
orders are linear on each piece.  Divisorial rational rays are dense
in each rational piece.  The exact zero assertions extend by continuity;
applying this separately to each function gives triviality on the
whole field.
\end{proof}

\subsection{Adjoint lattice points for integral Weil polarizations}
\label{subsec:C-adjoint}

We prove the precise lattice-point supply needed for the lower count.
The polarization will be nef and big and may be an integral Weil
divisor with unbounded Cartier index.  The rounding and pushdown steps
below keep those hypotheses throughout.

The weighted-blowup and adjoint-lifting construction below adapts the
``Adjoint lattice points'' and ``Orbifold vanishing'' lemmas in
\cite{OpenAIAdjointSemigroup}.  Here the argument is developed for a nef
and big integral Weil polarization with a klt Calabi--Yau boundary,
including the rounding and global regularity checks given below; the
earlier theorem is not used as a black box.

\begin{lemma}[Vanishing on a smooth orbifold]\label{lem:C-orbifold-vanishing}
Let $\mathcal X$ be a smooth proper Deligne--Mumford stack over an
algebraically closed field of characteristic zero, with trivial generic
stabilizer and projective coarse space.  Let $\Theta$ be an effective
rational divisor with simple normal crossing support and coefficients
less than one.  If $D$ is an integral Cartier divisor and
$D-(K_{\mathcal X}+\Theta)$ is nef and big, then
$H^q(\mathcal X,\mathcal O_{\mathcal X}(D))=0$ for $q>0$.
\end{lemma}
\begin{proof}
Write $q:\mathcal X\to X$ for the coarse morphism.  For each prime $B$
of $X$, let $P$ be its reduced inverse-image divisor and $e_B$ its
generic stabilizer order.  Thus $q^*B=e_BP$.  Rational divisors on
$\mathcal X$ descend uniquely to rational divisors on $X$; use a
subscript $c$ for the descended divisor.  The local normal-direction
equation $t=z^{e_B}$ gives
\[
 K_{\mathcal X}=q^*(K_X+R),\qquad
 R=\sum_B(1-e_B^{-1})B.
\]
Choose compatible canonical divisors using one rational top form.
Set $\Delta=R+\Theta_c-\{D_c\}$.  If the coefficients of $D$ and
$\Theta$ at $P$ are $d\in\Z$ and $\theta\in[0,1)$, write
$d=e_Bn+s$ with $0\le s<e_B$, also when $d$ is negative.  Then
\[
 \operatorname{coeff}_B\Delta
    =\frac{e_B-1+\theta-s}{e_B}\in[0,1).
\]
In particular $\Delta$ is effective.

The pair $(X,\Delta)$ is klt, including at higher-codimensional quotient
singularities.  Indeed, etale locally on $X$ there is a finite quotient
presentation $[U/G]$, with $U$ smooth and $f:U\to U/G$ finite.  Its
log pullback boundary is
$\Theta_U-f^*\{D_c\}$.  This possibly negative boundary is bounded
above by the simple normal crossing boundary $\Theta_U$.  Thus for
every divisorial valuation $F$ over $U$,
\[
 A(F;U,\Theta_U-f^*\{D_c\})
   =A(F;U,\Theta_U)+F(f^*\{D_c\})>0.
\]
The finite-map discrepancy formula multiplies downstairs log
discrepancies by positive ramification indices, so the downstairs
pair is klt.  This argument requires only local quotient charts.

The coarse pushforward is
\[
 q_*\mathcal O_{\mathcal X}(D)=\mathcal O_X(\lfloor D_c\rfloor).
\]
To check it, a rational section $h$ is regular along $P$ precisely when
$e_B\operatorname{ord}_B(h)+d\ge0$, equivalently
$\operatorname{ord}_B(h)+\lfloor d/e_B\rfloor\ge0$.  Both sides
extend these codimension-one conditions across codimension at least two:
on a smooth finite quotient chart this follows by normal extension and
then taking invariants.  The stack is tame in characteristic zero, so
coarse pushforward is exact and identifies the cohomology groups.

Quotient singularities are locally $\Q$-factorial.  Consequently
$M=\lfloor D_c\rfloor$ is an integral $\Q$-Cartier Weil divisor,
and
\[
 M-(K_X+\Delta)
   =D_c-(K_X+R+\Theta_c)
\]
is nef and big.  Kawamata--Viehweg vanishing for a klt pair and an
integral $\Q$-Cartier Weil divisor gives the result.  We use \cite[Theorem~3.2.4, manuscript version of 9 October 2023]{FujinoVanishing};
the additional log-bigness condition on log canonical centres is
vacuous for a klt pair.  The characteristic-zero statement also follows
from its complex form by descending the finite-type data to a finitely
generated field and using invariance of nefness, bigness and cohomology
under algebraically closed field extension.
\end{proof}

\begin{lemma}[Adjoint Taylor exponents]\label{lem:C-adjoint-lattice}
Let $X$ be a projective $\Q$-factorial variety of Fano type of
dimension $p>0$, and let $N$ be a nef and big integral Weil divisor.
Suppose that $H^0(X,\mathcal O_X(N))$ defines a birational map.
Choose an effective rational boundary $\Delta$ with $(X,\Delta)$ klt
and $K_X+\Delta\sim_{\Q}0$.  At a general smooth point outside
the supports of $N$ and $\Delta$, fix regular parameters and take
total-degree-then-lexicographic Taylor initial exponents, using the
compatible local trivializations of the powers of $N$.  Let
\[
 K=\overline{\operatorname{conv}}
 \left\{\frac{\nu(s)}j:
 j\in\Z_{>0},\quad
 0\ne s\in H^0(X,\mathcal O_X(jN))\right\}.
\]
If $t$ is a positive integer and $\alpha\in\Z_{\ge0}^p$ satisfies
$\alpha+\mathbf1\in\operatorname{int}(tK)$, then some section of
$tN$ has initial exponent $\alpha$.
\end{lemma}
\begin{proof}
The shift by $\mathbf1$ supplies the adjoint correction in this
argument.  A weighted blowup will convert it into the canonical
discrepancy, and the same shift will place the required monomial in a
multiplier ideal on the exceptional divisor.  Vanishing then lifts
that monomial to a global section.

\medskip
\noindent\emph{Preparing sections with weighted initials.}
Resolve $(X,\Delta)$ by $f:Y\to X$, isomorphically over a neighborhood
of the chosen point $x$.  Put $N_Y=f^*N$ and define the log pullback
$\Delta_Y$ by
\[
 K_Y+\Delta_Y=f^*(K_X+\Delta).
\]
The exceptional coefficients of $\Delta_Y$ may be negative.  This
will cause no difficulty: our vanishing boundary will be a fractional
rounding correction, not $\Delta_Y$.  Both $N_Y$ and $\Delta_Y$ have
zero divisor support near $x$ after compatible local trivialization.
Choose canonical divisors with no support near $x$ as well.

Let $C$ be the closed cone generated by the actual pairs $(j,\nu(s))$.
It has slice $K$ at degree one.  Choose rational $0<t_0<t$ with
$(t_0,\alpha+\mathbf1)\in\operatorname{int}C$.  There is a cone
generated by finitely many actual pairs that already contains this
point in its interior: approximate the vertices of a small simplex
around the point by finite nonnegative combinations of actual pairs.
Sufficiently close approximations still surround the point.

Choose positive primitive integral weights $u_1,\ldots,u_p$ that
replace the Taylor order by scalar weight on this finite list, each
listed section having its specified initial as a single weight initial.
Also require
\[
 \gamma\text{ precedes }\alpha\quad\Longrightarrow\quad
        u\cdot\gamma<u\cdot\alpha.
\]
Here is the finite-order justification.  Take an integer $d$ at least
the total degree of every listed initial and of $\alpha$.  Weights
sufficiently close to $(1,\ldots,1)$ preserve every strict total-degree
comparison through degree $d$, and satisfy
$(d+1)\min u_i>d\max u_i$, so all higher-degree monomials are larger
than every relevant initial.  Successive sufficiently small rational
perturbations realize lexicographic comparisons within the finitely
many monomials of degree at most $d$.  Clearing denominators and
dividing by the common divisor preserves all strict comparisons.
This also proves the assertion for formal Taylor expansions.

Put $p_0=u\cdot\alpha$ and $U=u\cdot(\alpha+\mathbf1)$.
The finite cone meets the plane with degree $t_0$ and weight $U$ in
a neighborhood of $(t_0,\alpha+\mathbf1)$ relative to that entire
plane.  Choose finitely many rational points surrounding it there.
Write these points as nonnegative rational combinations of the cone
generators and clear denominators.  Taking products, and then common
powers as necessary, gives a sufficiently divisible positive integer
$l$ and sections of the Cartier divisor $lt_0N_Y$ with single weight
initials $z^{\beta^{(1)}},\ldots,z^{\beta^{(r)}}$ such that
\[
 u\cdot\beta^{(a)}=lU,\qquad
 \alpha+\mathbf1\in
 \operatorname{relint}\operatorname{conv}
       \{\beta^{(1)}/l,\ldots,\beta^{(r)}/l\}.
\]
The relative interior contains a neighborhood in the whole plane
$u\cdot v=U$.  Pulling these sections to $Y$ is legitimate after
clearing denominators: the effective divisor of each product is
$\Q$-Cartier and its rational pullback is effective.

\medskip
\noindent\emph{The weighted modification and the vanishing divisor.}
Let $\pi:\mathcal Z\to Y$ be the smooth stack weighted blowup at $x$
with weights $u$, and let $E$ be its exceptional Cartier divisor.
On the coordinate model it is
\[
 \left[(\mathbb A^{p+1}_{x_1,\ldots,x_p,\tau}
       \setminus\{x_1=\cdots=x_p=0\})/\mathbb G_m\right],
 \qquad
 \rho\cdot(x_i,\tau)=(\rho^{u_i}x_i,\rho^{-1}\tau),
 \qquad z_i=\tau^{u_i}x_i.
\]
Its exceptional divisor is the weighted projective stack $\mathbb P(u)$.
This description pulls back along an etale parameter map near $x$
and glues to the identity away from $x$.  Equivalently use the weighted
ideal filtration generated by monomials of weight at least $k$; each
ideal is supported on a finite infinitesimal neighborhood of $x$ and
extends as the unit ideal outside $x$.  The finitely generated weighted
Rees algebra makes the coarse modification projective.  The stack is
smooth with finite stabilizers and trivial generic stabilizer.  On the
displayed atlas the Jacobian gives
\[
 K_{\mathcal Z}=\pi^*K_Y+(\textstyle\sum_i u_i-1)E.
\]
For $p=1$ this construction is the identity modification and $E=x$;
the same formulas and restriction sequence apply.

Write $\Delta_{\mathcal Z}=\pi^*\Delta_Y$, supported away from $E$,
and define
\[
 P_0=\pi^*(tN_Y)-p_0E,\qquad
 H_0=\pi^*(lt_0N_Y)-lUE.
\]
The divisor $P_0$ is rational globally but integral Cartier near $E$;
$H_0$ is Cartier globally.  The selected sections, after their common
weighted order is removed, are sections of $H_0$.  Their base ideal
$\mathfrak b$ does not vanish identically on $E$, and its restriction
to $E$ is generated by their monomial initials.  Take a projective
log resolution $\sigma:\mathcal Y\to\mathcal Z$ of these data and
all rational divisor supports used below, with strict transform
$\widetilde E$.  Arrange that the total support is simple normal
crossing, that $\sigma$ is an isomorphism at the generic point of
$E$, and that $\tau:\widetilde E\to E$ resolves $\mathfrak b|_E$.
First resolve the additional divisor supports, then principalize the
ideal with the resulting simple normal crossing boundary, including
the transform of $E$, on smooth charts
\cite[Theorems~2.4.1--2.4.3]{WlodarczykResolution}.  Smooth functoriality
identifies these constructions on chart overlaps, so they descend to
the required smooth stack; the sequence of blowups has projective
coarse space.  Write
\[
 \mathfrak b\mathcal O_{\mathcal Y}=\mathcal O_{\mathcal Y}(-G).
\]
Then $G\ge0$ does not contain $\widetilde E$, and
$\sigma^*H_0-G$ is globally generated.

Set
\[
 Q=P_0-K_{\mathcal Z}-\Delta_{\mathcal Z}-E,
 \quad R=\sigma^*Q-G/l,
 \quad A=K_{\mathcal Y}+\lceil R\rceil,
 \quad\Theta=\lceil R\rceil-R.
\]
Thus $A$ is an integral Cartier divisor, and $\Theta$ is effective
with simple normal crossing support and coefficients in $[0,1)$.
The weighted canonical identity, $U-p_0=\sum u_i$, and
$K_Y+\Delta_Y\sim_{\Q}0$ give
\[
 Q-H_0/l\sim_{\Q}(t-t_0)\pi^*N_Y.
\]
Consequently
\[
 A-(K_{\mathcal Y}+\Theta)=R
 \sim_{\Q}
 (t-t_0)\sigma^*\pi^*N_Y+(\sigma^*H_0-G)/l.
\]
The first summand is nef and big and the second is nef.  Lemma~\ref{lem:C-orbifold-vanishing} therefore gives $H^1(\mathcal Y,\mathcal O(A))=0$ and
the restriction surjection
\[
 H^0(\mathcal Y,A+\widetilde E)
 \longrightarrow
 H^0(\widetilde E,(A+\widetilde E)|_{\widetilde E}).
\]

\medskip
\noindent\emph{The boundary monomial.}
We now exhibit the required section on the exceptional divisor.
Near $E$ the original
boundary vanishes and $P_0$ is integral.  Thus, on a neighborhood of
$\widetilde E$,
\[
 A=K_{\mathcal Y}
   +\sigma^*(P_0-K_{\mathcal Z}-E)-\lfloor G/l\rfloor.
\]
Adjunction identifies the restricted bundle with
\[
 \tau^*(P_0|_E)+K_{\widetilde E}-\tau^*K_E
                   -\lfloor(G|_{\widetilde E})/l\rfloor.
\]
Round-down commutes with this restriction: simple normal crossings
prevent two components of $G$ from restricting to a common prime
divisor on $\widetilde E$, and each disconnected intersection retains
the same coefficient.  This is the log-resolution expression for
$P_0|_E\otimes\mathcal J_E((\mathfrak b|_E)^{1/l})$.

The weighted-homogeneous monomial $z^\alpha$ is a global section of
$P_0|_E=\mathcal O_E(p_0)$, after the fixed local frame.  On the chart
$x_k\ne0$, use the smooth finite atlas setting $x_k=1$, with residual
group $\mu_{u_k}$.  The ideal becomes the ordinary monomial ideal
generated by $x_{\mathrm{rem}}^{\beta^{(a)}_{\mathrm{rem}}}$ and the
proposed section becomes $x_{\mathrm{rem}}^{\alpha_{\mathrm{rem}}}$.
Dropping coordinate $k$ is an affine isomorphism from the plane
$u\cdot v=U$ to $\R^{p-1}$, because $u_k>0$.  Hence
$\alpha_{\mathrm{rem}}+\mathbf1$ lies in the ordinary interior of
the convex hull of the projected vectors $\beta^{(a)}_{\mathrm{rem}}/l$,
and therefore in the interior of the scaled Newton polyhedron.
Howald's monomial multiplier-ideal formula \cite[Main Theorem]{Howald}
gives the required membership
on this smooth atlas.  It is equivariant for $\mu_{u_k}$, so descends
on every chart.  No fractional coordinate exponents are introduced:
the root needed to set $x_k=1$ is already part of the atlas.

\medskip
\noindent\emph{Lifting and pushing down.}
The restriction surjection now lifts this monomial.  Near $E$ its
lifting bundle differs from $\sigma^*P_0$ by
\[
 K_{\mathcal Y}-\sigma^*K_{\mathcal Z}
  -\sigma^*E+\widetilde E-\lfloor G/l\rfloor.
\]
The part preceding the last term is exceptional over $\mathcal Z$;
indeed it is effective by the smooth blowup discrepancy calculation
for the smooth pair $(\mathcal Z,E)$.  At every nonexceptional prime
the displayed correction is nonpositive.  Therefore the lifted
rational section pushes down to a regular section of $P_0$ near $E$.
Its restriction to $E$ is $z^\alpha$, initially generically and then
everywhere by regularity.  Pushing further to the smooth neighborhood
of $x$ gives a regular section with weight initial exactly $z^\alpha$.

It remains to check global regularity downstairs, since $N_Y$ was
rational and $\Delta_Y$ could have negative coefficients.  At an
original prime $B$ of $X$ away from $x$, write
$n_B=\operatorname{coeff}_B(tN)\in\Z$,
$\delta_B=\operatorname{coeff}_B\Delta\ge0$, and let $g_B\ge0$ be
the corresponding coefficient of $G$.  Canonical coefficients cancel
inside the ceiling, and the allowed divisor coefficient of the
pushed lifting bundle is
\[
 n_B+\lceil-\delta_B-g_B/l\rceil\le n_B.
\]
Thus no forbidden pole appears at any original prime.  The point
$x$ was already handled by the local weighted calculation (also when
$p=1$).  Normality now gives a section of $\mathcal O_X(tN)$.
Its weight initial excludes every monomial preceding $\alpha$ in
the original Taylor order, by the choice of $u$, and the coefficient
of $z^\alpha$ is nonzero.  Its Taylor initial is therefore $\alpha$.
\end{proof}

\begin{lemma}[Uniform count from the jet body]\label{lem:C-jet-count}
In the setting of Lemma~\ref{lem:C-adjoint-lattice}, $K$ is bounded and contains the standard
simplex $\Sigma=\operatorname{conv}(0,e_1,\ldots,e_p)$.  For every
integer $n\ge1$,
\[
 h^0(X,nN)\le(n+p)^p\operatorname{vol}(K).
\]
Fix any integer $t>1+2p$.  The space spanned by products of $n$ sections
of $tN$ has dimension at least $n^p\operatorname{vol}(K)$.
\end{lemma}
\begin{proof}
On a nonempty open set, the map given by $|N|$ is an immersion and
has no base point, because it is birational in characteristic zero.
At the chosen general point its values and first derivatives thus
span all first jets.  Suitable linear combinations give initial
exponents $0,e_1,\ldots,e_p$, proving $\Sigma\subset K$.
The order of vanishing at this point of a section of $jN$ is at most
$Cj$: intersect its effective divisor with a general complete
intersection curve through the point, chosen not to lie in its
support.  Its degree is $jN\cdot H^{p-1}$ for a fixed sufficiently
ample $H$, and bounds the local intersection multiplicity.  The
constant may depend on this individual example, which suffices to
make $K$ bounded.

For a finite-dimensional space of germs, the number of distinct
leading exponents is its dimension, by successive elimination of
leading coefficients; alternatively use the filtration by the
well-ordered monomials, whose nonzero graded pieces are one-dimensional.
The exponent set at degree $n$ is contained in $nK\cap\Z^p$.
The disjoint half-open unit boxes based at these lattice points lie
in
\[
 nK+[0,1]^p\subset nK+p\Sigma\subset(n+p)K.
\]
Their total volume is their number, proving the upper bound.

For the lower bound let $y=nx\in nK$.  Write $x_i=a_i+f_i$ with
$a_i\in\Z_{\ge0}$ and $0\le f_i<1$.  Choose $n$ integers
in coordinate $i$, namely $a_i+1$ in exactly $\lfloor nf_i\rfloor$
of the positions and $a_i$ in the others.  Doing this independently
in every coordinate gives $\alpha^{(1)},\ldots,\alpha^{(n)}$ with
\[
 \sum_{r=1}^n\alpha^{(r)}=\lfloor y\rfloor,
 \qquad 0<\alpha_i^{(r)}+1-x_i<2.
\]
When $f_i=0$ all entries equal $a_i$, so the strict inequalities still
hold.  Therefore
$\alpha^{(r)}+\mathbf1-x\in\operatorname{int}((t-1)\Sigma)$:
all coordinates are positive and their sum is less than $2p<t-1$.
Since $\Sigma\subset K$ and $x\in K$, a neighborhood of
$\alpha^{(r)}+\mathbf1$ lies in
$K+(t-1)K=tK$.  Lemma~\ref{lem:C-adjoint-lattice} supplies a degree-$t$ section
with each exponent $\alpha^{(r)}$.  Their product has exponent
$\lfloor y\rfloor$.

Consequently all lattice points $\lfloor y\rfloor$, for $y\in nK$,
occur among product initials.  Their unit boxes cover $nK$, so their
number is at least $\operatorname{vol}(nK)=n^p\operatorname{vol}(K)$.
Sections with distinct initials are linearly independent, proving
the claimed lower bound.  The case $p=0$ is immediate separately.
\end{proof}

\subsection{Completion of the saturation count}

The field and exact-triviality assertions are already established.
We finish by placing both section counts on one nef integral Weil
polarization.  The preceding two lemmas then give the same jet-body
volume as the normalization in the upper and lower estimates.
\label{subsec:C-count-completion}

\begin{proof}[Proof of part \textup{(iii)} of Lemma~\ref{lem:C-saturation-count}]
If $p=0$, then $K_*=k$.  By part (i), each fixed-cost section space
is eventually the one-dimensional constant space.  Take $v_i=1$.
Assume $p>0$ and use the extraction $\widehat W$ from
Lemma~\ref{lem:C-cutoff-claims} at $\beta_1$ of
\eqref{eq:C-plateau}.  Let $D_*$ be the pullback of $D_W$ to
$\widehat W$.  For every rational $m>0$, the cutoff spaces identify
exactly as function spaces:
\begin{equation}\label{eq:C-cutoff-base-spaces}
                       S(mP_{\beta_1})=S(mD_*).
\end{equation}
The negative-MMP error is effective and exceptional, so pushing and
pulling the divisor inequality for an individual rational function
identifies the space with the one on the semiample output.  A function
in that space belongs to $k(W)$; the equality $P_{\beta_1}'=p^*D_W$
then identifies its inequalities with those downstairs.  Pullback to
$\widehat W$ does not change them.  This argument works for each
rational $m$, with no requirement that $mP_{\beta_1}$ be Cartier.

Every positive coefficient of $LD_*$ is at least $c'\delta$, because
$L\beta_1\geq c'$ and the coefficients of $D_*/\beta_1$ are at
least $\delta$.  Choose a fixed sufficiently large integer $M$ and put
\begin{equation}\label{eq:C-count-polarization}
                              N=\lfloor MLD_*\rfloor.
\end{equation}
The coefficient lower bound gives $N\geq LD_*$ once $M$ is large
enough.  Increase $M$ so that $N\geq C D_*/\beta_1$, with the
constant from the field-generation construction.  Then $S(N)$ defines
a birational map and $N$ is big and integral.  Run its MMP to a nef
and big integral Weil divisor $N'$ on a $\Q$-factorial projective
Fano type variety $X'$.  For every positive integer $j$, negativity
gives
\begin{equation}\label{eq:C-count-mmp-spaces}
                     S(jN)=S(jN'),\qquad
                              S(jN')\subset S(jMLD).
\end{equation}
Indeed the first equality follows by pushing an effective divisor and,
in the reverse direction, adding the effective exceptional MMP error;
the second follows from $N\leq MLD_*$ and
\eqref{eq:C-cutoff-base-spaces}.  For every fixed positive integer $n$,
part (i) and the exact plateau identities give eventually
\begin{equation}\label{eq:C-count-upper-inclusion}
               S(nLD)\subset S(nLD_*)\subset S(nN)=S(nN').
\end{equation}

Choose an effective ordinary rational klt Calabi--Yau boundary
$\Delta$ on $X'$.  Such a boundary is obtained from a klt Fano type
boundary by adding a general sufficiently divisible rational member
of its ample anti-log-canonical class.  At a general smooth point
outside the boundary and divisor supports, let $K_i$ be the jet body
of $N'$ in Lemma~\ref{lem:C-adjoint-lattice}.  Its volume is positive
and finite by Lemma~\ref{lem:C-jet-count}.  Set
\[
                              v_i=\vol(K_i).
\]
Equations \eqref{eq:C-count-upper-inclusion} and
Lemma~\ref{lem:C-jet-count} imply
\[
 \dim S(nLD)\leq(n+p)^p v_i
                    \leq(p+1)^p v_i n^p.
\]
Fix once and for all an integer $t>1+2p$.  The same lemma supplies a
space of products of $n$ sections of $tN'$ of dimension at least
$v_i n^p$.  By \eqref{eq:C-count-mmp-spaces},
$S(tN')\subset S(tMLD)$.  Enlarge the constant $B$ in part (ii) to
at least $tM$.  Those products then lie in the product space in
\eqref{eq:C-uniform-counts}, proving the lower bound with the same
$v_i$.  Both $t$ and $M$, and hence $B$, were fixed before $n$.
This completes the lemma.
\end{proof}

\begin{corollary}[Fixed generators and neutral algebraic extensions]
\label{cor:C-neutral-field}
In Lemma~\ref{lem:C-saturation-count}, let $g_{i,1},\ldots,g_{i,r_i}$
be any finite list contained in $S(C_1L_iD_i)$ for one fixed $C_1$.
For all sufficiently large $i$, the entire field
$F_i=k(g_{i,1},\ldots,g_{i,r_i})$, and its relative algebraic closure
inside $k(V_i)$, lie in $K_{*,i}$.  Thus every test from part
\textup{(ii)} is exactly zero on every nonzero element of these fields
at that same index, without a degree-dependent eventual qualification.
More generally, a valuation trivial on a subfield remains trivial on
any algebraic extension of it inside its valued overfield.
\end{corollary}

\begin{proof}
Part (i) of the lemma includes the whole space $S(C_1L_iD_i)$ at a
single eventual index; it does not supply separate indices for its
individual elements.  Field closure gives $F_i\subset K_{*,i}$.
The cutoff field is relatively algebraically closed, so it contains
the stated relative algebraic closure as well.  To see the last
assertion directly, let $x\ne0$ be algebraic and choose a polynomial
relation for it with nonzero constant term and coefficients in the
trivially valued subfield.  If the value of $x$ were positive, the
constant term would have uniquely least value; if it were negative,
the highest nonzero term would have uniquely least value.  Neither is
possible in a vanishing sum.  Hence its value is zero.  This applies
to every rational expression and every polynomial degree once the
finite generators have been included, rather than by a further limit
in that degree.
\end{proof}
